\frfilename{mt135.tex}
\versiondate{14.9.04/14.7.07}
\copyrightdate{1994}

\def\chaptername{Pelengkap}
\def\sectionname{Garis real diperluas}

\newsection{135}

Sering kali ada gunanya mengizinkan
`$\infty$' masuk ke dalam rumus-rumus kita, dan dalam konteks teori ukuran
manipulasi yang tepat cukup konsisten sehingga memungkinkan kita
mengembangkan teori tentang {\bf garis real diperluas}, yaitu himpunan
$[-\infty,\infty]=\Bbb R\cup\{-\infty,\infty\}$, yang kadang ditulis
$\overline{\Bbb R}$. Saya memberikan uraian singkat tanpa bukti lengkap,
karena saya berharap bahwa ketika materi ini mulai diperlukan dalam
argumen-argumen yang saya gunakan, semuanya akan tampak sepenuhnya elementer.

\leader{135A}{Struktur aljabar dari $[-\infty,\infty]$ (a)} Jika kita
menuliskan

\Centerline{$a+\infty=\infty+a=\infty$,
\quad$a+(-\infty)=(-\infty)+a=-\infty$}

\noindent untuk setiap $a\in\Bbb R$, dan

\Centerline{$\infty+\infty=\infty$,\quad$(-\infty)+(-\infty)=-\infty$,}

\noindent tetapi tidak mendefinisikan $\infty+(-\infty)$ atau
$(-\infty)+\infty$, kita memperoleh suatu operasi biner yang terdefinisi
secara parsial pada $[-\infty,\infty]$, yang memperluas penjumlahan biasa
pada $\Bbb R$. Operasi ini bersifat {\it asosiatif} dalam arti bahwa

\inset{jika $u$, $v$, $w\in[-\infty,\infty]$ dan salah satu dari
$u+(v+w)$, $(u+v)+w$ terdefinisi, maka yang lainnya juga terdefinisi,
dan keduanya sama,}

\noindent dan bersifat {\it komutatif} dalam arti bahwa

\inset{jika $u$, $v\in[-\infty,\infty]$ dan salah satu dari $u+v$, $v+u$
terdefinisi, maka yang lainnya juga terdefinisi, dan keduanya sama.}

\noindent Operasi ini mempunyai unsur {\it identitas} $0$ sedemikian
sehingga $u+0=0+u=u$ untuk setiap $u\in[-\infty,\infty]$; tetapi $\infty$
dan $-\infty$ tidak memiliki invers.

\header{135Ab}{\bf (b)} Jika kita mendefinisikan

\Centerline{$a\cdot\infty=\infty\cdot a=\infty$,
\quad$a\cdot(-\infty)=(-\infty)\cdot a=-\infty$}

\noindent untuk bilangan real $a>0$,

\Centerline{$a\cdot\infty=\infty\cdot a=-\infty$,
\quad$a\cdot(-\infty)=(-\infty)\cdot a=\infty$}

\noindent untuk bilangan real $a<0$,

\Centerline{$\infty\cdot\infty=(-\infty)\cdot(-\infty)=\infty$,
\quad$(-\infty)\cdot\infty=\infty\cdot(-\infty)=-\infty$,}

\Centerline{$0\cdot\infty=\infty\cdot 0=0\cdot(-\infty)
=(-\infty)\cdot0=0$}

\noindent maka kita memperoleh suatu operasi biner pada
$[-\infty,\infty]$ yang memperluas perkalian biasa pada $\Bbb R$, yang
bersifat asosiatif dan komutatif serta mempunyai identitas 1; $0$, $\infty$
dan $-\infty$ tidak memiliki invers.

\header{135Ac}{\bf (c)} Kita mempunyai suatu {\it hukum distributif}
\cmmnt{, yang sedikit lebih lemah daripada hukum asosiatif dan komutatif
untuk penjumlahan}:

\inset{jika $u$, $v$, $w\in[-\infty,\infty]$ dan baik $u(v+w)$ maupun
$uv+uw$ terdefinisi, maka keduanya sama.}

\cmmnt{\noindent (Namun perhatikan masalah yang timbul pada kombinasi
seperti $\infty(1+(-2))$, $0\cdot\infty+0\cdot(-\infty)$.)}

\header{135Ad}{\bf (d)} Meskipun $\infty$ dan $-\infty$ tidak mempunyai
invers dalam semigrup $([-\infty,\infty],\cdot)$, tampaknya tidak ada
masalah jika kita menuliskan $a/\infty=a/(-\infty)=0$ untuk setiap
$a\in\Bbb R$.
\cmmnt{Namun tentu saja perluasan gagasan pembagian seperti ini harus
diperlakukan dengan hati-hati dalam rumus seperti $u\cdot\bover{v}{u}$.}

\leader{135B}{Struktur urutan dari $[-\infty,\infty]$ (a)} Jika kita
menuliskan

\Centerline{$-\infty\le u\le \infty$ untuk setiap
$u\in[-\infty,\infty]$,}

\noindent kita memperoleh suatu relasi pada $[-\infty,\infty]$, yang
memperluas urutan biasa pada $\Bbb R$, dan merupakan urutan {\it total},
yaitu,

\inset{untuk sembarang $u$, $v$, $w\in[-\infty,\infty]$, jika $u\le v$
dan $v\le w$ maka $u\le w$,}

\inset{$u\le u$ untuk setiap $u\in[-\infty,\infty]$,}

\inset{untuk sembarang $u$, $v\in[-\infty,\infty]$, jika $u\le v$ dan
$v\le u$ maka $u=v$,}

\inset{untuk sembarang $u$, $v\in[-\infty,\infty]$, berlaku $u\le v$
atau $v\le u$.}

\noindent Selain itu, setiap subhimpunan dari $[-\infty,\infty]$ mempunyai
supremum dan infimum, jika kita menuliskan $\sup\emptyset=-\infty$,
$\inf\emptyset=\infty$.

\header{135Bb}{\bf (b)} Urutan ini `invarian terhadap translasi' dalam
arti lemah bahwa

\inset{jika $u$, $v$, $w\in[-\infty,\infty]$ dan $v\le w$, serta $u+v$,
$u+w$ keduanya terdefinisi, maka $u+v\le u+w$.}

\noindent Urutan ini dipertahankan oleh perkalian dengan faktor nonnegatif dalam arti
bahwa

\inset{jika $u$, $v$, $w\in[-\infty,\infty]$ dan $0\le u$ serta
$v\le w$, maka $uv\le uw$,}

\noindent sedangkan urutan ini dibalik oleh perkalian dengan faktor nonpositif dalam
arti bahwa

\inset{jika $u$, $v$, $w\in[-\infty,\infty]$ dan $u\le 0$ serta
$v\le w$, maka $uw\le uv$.}

\leader{135C}{Struktur Borel dari $[-\infty,\infty]$} Kita mengatakan
bahwa suatu himpunan $E\subseteq[-\infty,\infty]$ adalah {\bf himpunan
Borel} dalam $[-\infty,\infty]$ jika $E\cap\Bbb R$ merupakan subhimpunan
Borel dari $\Bbb R$. Mudah diperiksa bahwa keluarga himpunan-himpunan
seperti ini merupakan suatu aljabar-$\sigma$ atas subhimpunan-subhimpunan
$[-\infty,\infty]$. \cmmnt{Lihat juga 135Xb di bawah.}

\ifnum\stylenumber=11\ifresultsonly\eject\fi\fi

\leader{135D}{Barisan konvergen dalam $[-\infty,\infty]$} Kita dapat
mengatakan bahwa suatu barisan $\sequencen{u_n}$ dalam
$[-\infty,\infty]$ {\bf konvergen} ke $u\in[-\infty,\infty]$ jika

\inset{kapan pun $v<u$ terdapat $n_0\in\Bbb N$ sedemikian sehingga
$v\le u_n$ untuk setiap $n\ge n_0$, dan kapan pun $u<v$ terdapat
$n_0\in\Bbb N$ sedemikian sehingga $u_n\le v$ untuk setiap $n\ge n_0$;}

\noindent atau, secara ekuivalen,

\inset{{\it entah} $u\in\Bbb R$ dan untuk setiap $\delta>0$
terdapat $n_0\in\Bbb N$ sedemikian sehingga
$u_n\in[u-\delta,u+\delta]$ untuk setiap $n\ge n_0$}

\inset{{\it atau} $u=-\infty$ dan untuk setiap $a\in\Bbb R$ terdapat
$n_0\in\Bbb N$ sedemikian sehingga $u_n\le a$ untuk setiap $n\ge n_0$}

\inset{{\it atau} $u=\infty$ dan untuk setiap $a\in\Bbb R$ terdapat
$n_0\in\Bbb N$ sedemikian sehingga $u_n\ge a$ untuk setiap $n\ge n_0$.}

\cmmnt{\noindent (Bandingkan dengan gagasan konvergensi dalam 112Ba.)}

\vleader{48pt}{135E}{Fungsi terukur} Misalkan $X$ sembarang himpunan dan
$\Sigma$ suatu aljabar-$\sigma$ atas subhimpunan-subhimpunan $X$.

\header{135Ea}{\bf (a)} Misalkan $D$ suatu subhimpunan dari $X$ dan $\Sigma_D$
adalah aljabar-$\sigma$ subruang (121A). Untuk sembarang fungsi
$f:D\to[-\infty,\infty]$, pernyataan-pernyataan berikut ekuivalen:

\quad (i) $\{x:f(x)<u\}\in\Sigma_D$ untuk setiap
$u\in[-\infty,\infty]$;

\quad (ii) $\{x:f(x)\le u\}\in\Sigma_D$ untuk setiap
$u\in[-\infty,\infty]$;

\quad (iii) $\{x:f(x)>u\}\in\Sigma_D$ untuk setiap
$u\in[-\infty,\infty]$;

\quad (iv) $\{x:f(x)\ge u\}\in\Sigma_D$ untuk setiap
$u\in[-\infty,\infty]$;

\quad (v) $\{x:f(x)\le q\}\in\Sigma_D$ untuk setiap $q\in\Bbb Q$.

\prooflet{\noindent\Prf\ Buktinya hampir identik dengan bukti 121B.
Satu-satunya perubahan adalah:

-- dalam (i)$\Rightarrow$(ii), $\{x:f(x)\le\infty\}$ dan
$\{x:f(x)\le -\infty\}$ tidak mesti sama dengan
$\bigcap_{n\in\Bbb N}\{x:f(x)<\infty+2^{-n}\}$,
$\bigcap_{n\in\Bbb N}\{x:f(x)<-\infty+2^{-n}\}$; tetapi yang pertama
adalah $D$, sehingga tentu termasuk dalam $\Sigma_D$, dan yang kedua adalah
$\bigcap_{n\in\Bbb N}\{x:f(x)<-n\}$, sehingga termasuk dalam $\Sigma_D$.

-- Dalam (iii)$\Rightarrow$(iv), dengan cara serupa, kita harus menggunakan
fakta-fakta bahwa

\Centerline{$\{x:f(x)\ge -\infty\}=D\in\Sigma_D$,
\quad$\{x:f(x)\ge\infty\}=\bigcap_{n\in\Bbb N}\{x:f(x)>n\}\in\Sigma_D$.}

-- Mengenai syarat tambahan (v), tentu saja kita mempunyai
(ii)$\Rightarrow$(v), tetapi kita juga mempunyai (v)$\Rightarrow$(i), karena

\Centerline{$\{x:f(x)<u\}=\bigcup_{q\in\Bbb Q,q<u}\{x:f(x)\le q\}$}

\noindent untuk setiap $u\in[-\infty,\infty]$.
\Qed}


\header{135Eb}{\bf (b)} Karena itu kita dapat mengatakan\cmmnt{, seperti
dalam 121C,} bahwa suatu fungsi yang bernilai dalam $[-\infty,\infty]$
adalah {\bf terukur} jika fungsi tersebut memenuhi syarat-syarat ekuivalen
ini.

\header{135Ec}{\bf (c)} Perhatikan bahwa jika
$f:D\to[-\infty,\infty]$ bersifat $\Sigma$-terukur, maka

\Centerline{$E_{\infty}(f)=f^{-1}[\{\infty\}]=\{x:f(x)\ge\infty\}$,
\quad $E_{-\infty}(f)=f^{-1}[\{-\infty\}]=\{x:f(x)\le-\infty\}$}

\noindent harus termasuk dalam $\Sigma_D$, sedangkan
$f_{\Bbb R}=f\restr D\setminus(E_{\infty}(f)\cup E_{-\infty}(f))$, yaitu
`bagian bernilai real dari $f$', bersifat terukur dalam arti 121C.

\header{135Ed}{\bf (d)} Sebaliknya, jika $E_{\infty}$ dan $E_{-\infty}$
termasuk dalam $\Sigma_D$, dan
$f_{\Bbb R}:D\setminus(E_{\infty}\cup E_{-\infty})\to\Bbb R$ bersifat
terukur, maka $f:D\to[-\infty,\infty]$ akan bersifat terukur, dengan
$f(x)=\infty$ jika $x\in E_{\infty}$, $f(x)=-\infty$ jika
$x\in E_{-\infty}$ dan $f(x)=f_{\Bbb R}(x)$ untuk $x\in D$ lainnya.

\header{135Ee}{\bf (e)} Akibatnya, jika $f$, $g$ adalah fungsi-fungsi
terukur dari subhimpunan-subhimpunan $X$ ke $[-\infty,\infty]$, maka
$f+g$, $f\times g$ dan $f/g$ bersifat terukur. \prooflet{\Prf\ Hal ini
dapat dibuktikan dengan menyesuaikan argumen-argumen 121Eb, 121Ed dan 121Ee,
atau dengan menerapkan hasil-hasil tersebut pada $f_{\Bbb R}$ dan
$g_{\Bbb R}$ serta meninjau secara terpisah himpunan-himpunan tempat salah
satu atau keduanya bernilai tak berhingga. \Qed}

\header{135Ef}{\bf (f)} Kita dapat mengatakan bahwa suatu fungsi $h$ dari
subhimpunan $D$ dari $[-\infty,\infty]$ ke $[-\infty,\infty]$ adalah
{\bf terukur Borel} jika fungsi tersebut terukur\cmmnt{ (dalam arti (b) di
atas)} terhadap aljabar-$\sigma$ Borel dari $[-\infty,\infty]$
\cmmnt{ (sebagaimana didefinisikan dalam 135C)}. Sekarang, jika $X$ suatu
himpunan, $\Sigma$ suatu aljabar-$\sigma$ atas subhimpunan-subhimpunan $X$,
$f$ suatu fungsi terukur dari subhimpunan dari $X$ ke $[-\infty,\infty]$ dan
$h$ suatu fungsi terukur Borel dari subhimpunan dari $[-\infty,\infty]$ ke
$[-\infty,\infty]$, maka komposisi $hf$ terukur.
\prooflet{\Prf\ Terapkan 121Eg pada $h^*f_{\Bbb R}$, dengan
$h^*=h\restr(\Bbb R\cap h^{-1}[\Bbb R])$, lalu tinjau secara terpisah
himpunan-himpunan $\{x:f(x)=\pm\infty\}$,
$\{x:hf(x)=\pm\infty\}$. \Qed}

\header{135Eg}{\bf (g)} Misalkan $X$ suatu himpunan dan $\Sigma$ suatu
aljabar-$\sigma$ atas subhimpunan-subhimpunan $X$. Misalkan
$\sequencen{f_n}$ suatu barisan fungsi terukur dari
subhimpunan-subhimpunan $X$ ke $[-\infty,\infty]$. Maka
$\lim_{n\to\infty}f_n$, $\sup_{n\in\Bbb N}f_n$ dan
$\inf_{n\in\Bbb N}f_n$ bersifat terukur, jika\cmmnt{, mengikuti
prinsip-prinsip yang diuraikan dalam 121F,} kita menetapkan domain-domainnya
sebagai

\Centerline{$\{x:x\in\bigcup_{n\in\Bbb N}\bigcap_{m\ge n}\dom f_m,\,
\lim_{n\to\infty}f_n(x)$ ada dalam $[-\infty,\infty]\}$,}

\Centerline{$\bigcap_{n\in\Bbb N}\dom f_n$.}

\prooflet{\noindent\Prf\ Ikuti metode 121Fa-121Fc. \Qed}

\leader{135F}{Fungsi terintegralkan bernilai $[-\infty,\infty]$ (a)}
Tentu saja kita tidak akan menerima suatu fungsi sebagai `terintegralkan'
kecuali jika fungsi tersebut berhingga hampir di mana-mana, dan untuk
fungsi-fungsi seperti itu catatan-catatan dalam 133B sudah memadai.

\header{135Fb}{\bf (b)} Namun, kita dapat membuat perluasan yang konsisten
atas gagasan integral tak berhingga, dengan mengembangkan sedikit gagasan
dalam 133A. Jika $(X,\Sigma,\mu)$ suatu ruang ukur dan $f$ suatu fungsi,
yang terdefinisi hampir di mana-mana dalam $X$, bernilai dalam
$[0,\infty]$, dan terukur secara virtual (yaitu, sedemikian sehingga
$f\restr E$ terukur\cmmnt{ dalam arti 135E} untuk suatu himpunan
koterabaikan $E$), maka kita dapat dengan aman menuliskan
`$\int f=\infty$' kapan pun $f$ tidak terintegralkan. Kita akan mendapati
bahwa untuk fungsi-fungsi seperti itu berlaku
$\int f+g=\int f+\int g$ dan $\int cf=c\int f$ untuk setiap
$c\in[0,\infty]$, dengan menggunakan definisi penjumlahan dan perkalian
pada $[0,\infty]$ yang diberikan di atas. Akibatnya, seperti dalam
122M-122O, %122M 122N 122O
kita dapat mengatakan bahwa untuk suatu fungsi umum $f$ yang terukur secara
virtual, terdefinisi hampir di mana-mana dalam $X$, dan bernilai dalam
$[-\infty,\infty]$, berlaku $\int f=\int f_1-\int f_2$ kapan pun $f$
dapat dinyatakan sebagai selisih $f_1-f_2$ dari fungsi-fungsi nonnegatif
sedemikian sehingga $\int f_1$ dan $\int f_2$ keduanya terdefinisi dan
tidak keduanya tak berhingga. Sekarang kita mempunyai\cmmnt{, seperti
biasa,} rumus-rumus dasar

\Centerline{$\int f+g=\int f +\int g$,\quad
$\int cf=c\int f$,\quad
$\int|f|\ge|\int f|$}

\noindent kapan pun ruas-ruas kanan terdefinisi, dan $\int f\le\int g$
kapan pun $f\leae g$ dan kedua integral terdefinisi.
\cmmnt{Penting untuk diperhatikan bahwa} menurut definisi ini,
$\int f$ hanya dapat berhingga jika $f$ berhingga hampir di mana-mana.

\leader{135G}{}\cmmnt{ Sekarang kita mempunyai versi-versi teorema B. Levi
dan Lema Fatou (bandingkan 133K).

\wheader{135G}{6}{2}{2}{72pt}
\noindent}{\bf Proposisi} Misalkan $(X,\Sigma,\mu)$ suatu ruang ukur, dan
$\sequencen{f_n}$ suatu barisan fungsi bernilai $[-\infty,\infty]$ yang
terdefinisi hampir di mana-mana dalam $X$ dan mempunyai integral yang
terdefinisi dalam $[-\infty,\infty]$.

(a) Jika $f_n\leae f_{n+1}$ untuk setiap $n$ dan
$-\infty<\sup_{n\in\Bbb N}\int f_n$, maka
$\int\sup_{n\in\Bbb N}f_n=\sup_{n\in\Bbb N}\int f_n$.

(b) Jika, untuk setiap $n$, $f_n\ge 0$ hampir di mana-mana, maka
$\int\liminf_{n\to\infty}f_n\le\liminf_{n\to\infty}\int f_n$.

\proof{{\bf (a)} Perhatikan bahwa $f=\sup_{n\in\Bbb N}f_n$ terdefinisi
di setiap titik $\bigcap_{n\in\Bbb N}\dom f_n$, dan irisan ini
koterabaikan; selain itu, terdapat suatu himpunan koterabaikan $E$ sedemikian
sehingga $f_n\restr E$ terukur untuk setiap $n$, sehingga $f\restr E$
terukur. Sekarang, jika $u=\sup_{n\in\Bbb N}\int f_n$ berhingga, maka
semua kecuali sejumlah berhingga dari $f_n$ harus berhingga hampir di
mana-mana, dan hasilnya merupakan akibat teorema B. Levi untuk fungsi
bernilai real; sedangkan jika $u=\infty$ maka tentu
$\int\sup_{n\in\Bbb N}f_n$ tak berhingga.

\medskip

{\bf (b)} Seperti dalam 123B atau 133Kb, hasil ini sekarang diperoleh
dengan menerapkan (a) pada $g_n=\inf_{m\ge n}f_m$.
}%end of proof of 135G

\leader{135H}{Integral atas dan bawah sekali lagi (a)} Untuk menangani
fungsi yang bernilai dalam $[-\infty,\infty]$ kita perlu menyesuaikan
definisi-definisi dalam 133I. Misalkan $(X,\Sigma,\mu)$ suatu ruang ukur dan
$f$ suatu fungsi bernilai $[-\infty,\infty]$ yang terdefinisi hampir di
mana-mana dalam $X$. {\bf Integral atasnya} adalah

\Centerline{$\overlineint f
=\inf\{\int g:\int g$ terdefinisi dalam arti 135F dan $f\leae g\}$,}

\noindent dengan mengizinkan $\infty$ untuk $\inf\{\infty\}$ dan
$-\infty$ untuk $\inf\ocint{-\infty,\infty}$ atau
$\inf[-\infty,\infty]$. Dengan cara serupa, {\bf integral bawah} dari $f$
adalah

\Centerline{$\underlineint f
=\sup\{\int g:\int g$ terdefinisi, $f\geae g\}$.}

\noindent Dengan perubahan ini, semua hasil dalam 133J berlaku untuk
fungsi yang bernilai dalam $[-\infty,\infty]$, bukan hanya dalam $\Bbb R$.

\header{135Hb}{\bf (b)}\cmmnt{ Sejalan dengan 133Ka, kita mempunyai hasil
berikut.} Misalkan $(X,\Sigma,\mu)$ suatu ruang ukur, dan
$\sequencen{f_n}$ suatu barisan fungsi bernilai $[-\infty,\infty]$ yang
terdefinisi hampir di mana-mana dalam $X$.

\medskip

\quad{(i)} Jika $f_n\leae f_{n+1}$ untuk setiap $n$ dan
$\sup_{n\in\Bbb N}\overline{\int}f_n>-\infty$, maka
$\overline{\int}\sup_{n\in\Bbb N}f_n
=\sup_{n\in\Bbb N}\overline{\int}f_n$.

\medskip

\quad{(ii)} Jika, untuk setiap $n$, $f_n\ge 0$ hampir di mana-mana, maka
$\overline{\int}\liminf_{n\to\infty}f_n
\le\liminf_{n\to\infty}\overline{\int}f_n$.

\leader{135I}{\bf Ukuran \dvrocolon{subruang}}\cmmnt{ Kita perlu meninjau
kembali gagasan-gagasan dalam \S131 dalam konteks baru ini.

\medskip

\noindent}{\bf Proposisi} Misalkan $(X,\Sigma,\mu)$ suatu ruang ukur, dan
$H\in\Sigma$; tuliskan $\Sigma_H$ untuk aljabar-$\sigma$ subruang pada $H$
dan $\mu_H$ untuk ukuran subruang. Untuk sembarang fungsi bernilai
$[-\infty,\infty]$ yang kita lambangkan dengan $f$ dan terdefinisi pada suatu
subhimpunan dari $H$, tuliskan
$\tilde f$ untuk perluasan $f$ yang didefinisikan dengan menetapkan
$\tilde f(x)=f(x)$ jika $x\in\dom f$, dan $0$ jika
$x\in X\setminus H$.

(a) Andaikan $f$ suatu fungsi bernilai $[-\infty,\infty]$ yang terdefinisi
pada suatu subhimpunan dari $H$.

\quad(i) $\dom f$ bersifat $\mu_H$-koterabaikan jika dan hanya jika
$\dom\tilde f$ bersifat $\mu$-koterabaikan.

\quad(ii) $f$ bersifat $\mu_H$-terukur secara virtual jika dan hanya jika
$\tilde f$ bersifat $\mu$-terukur secara virtual.

\quad(iii) $\int_Hfd\mu_H=\int_X\tilde fd\mu$ jika salah satunya
terdefinisi dalam $[-\infty,\infty]$.

(b) Andaikan $h$ suatu fungsi bernilai $[-\infty,\infty]$ yang terdefinisi
hampir di mana-mana dalam $X$. Maka
$\int_H(h\restr H)d\mu_H=\int h\times\chi H\,d\mu$ jika salah satunya
terdefinisi dalam $[-\infty,\infty]$.

(c) Jika $h$ suatu fungsi bernilai $[-\infty,\infty]$ dan
$\int_Xh\,d\mu$ terdefinisi dalam $[-\infty,\infty]$, maka
$\int_H(h\restr H)d\mu_H$ terdefinisi dalam $[-\infty,\infty]$.

(d) Andaikan $h$ suatu fungsi bernilai $[-\infty,\infty]$ yang terdefinisi
hampir di mana-mana dalam $X$. Maka

\Centerline{$\overlineint_H(h\restr H)d\mu_H
=\overlineint_Xh\times\chi Hd\mu$.}

\proof{{\bf (a)(i)} Ini langsung diperoleh dari 131Ca, karena
$H\setminus\dom f=X\setminus\dom\tilde f$.

\medskip

\quad{\bf (ii)}\grheada\ Jika $f$ bersifat $\mu_H$-terukur secara
virtual, terdapat suatu himpunan $\mu_H$-koterabaikan $E\in\Sigma_H$
sedemikian sehingga $f\restr E$ bersifat $\Sigma_H$-terukur. Terdapat
$F\in\Sigma$ sedemikian sehingga $E=F\cap H$; sekarang
$G=F\cup(X\setminus H)$ termasuk dalam $\Sigma$, $E=G\cap H$, dan $G$
bersifat $\mu$-koterabaikan. Selain itu, untuk $q\in\Bbb Q$,

$$\eqalign{\{x:x\in G,\,\tilde f(x)\le q\}
&=\{x:x\in E,\,f(x)\le q\}\in\Sigma_H\subseteq\Sigma
   \text{ jika }q<0,\cr
&=\{x:x\in E,\,f(x)\le q\}\cup(X\setminus H)\in\Sigma
   \text{ jika }q\ge 0,\cr}$$

\noindent sehingga $\tilde f\restr G$ bersifat $\Sigma$-terukur dan
$\tilde f$ bersifat $\mu$-terukur secara virtual.

\medskip

\qquad\grheadb\ Jika $\tilde f$ bersifat $\mu$-terukur secara virtual,
terdapat suatu himpunan $\mu$-koterabaikan $G\in\Sigma$ sedemikian sehingga
$\tilde f\restr G$ bersifat $\Sigma$-terukur. Sekarang $E=G\cap H$
termasuk dalam $\Sigma_H$ dan bersifat $\mu_H$-koterabaikan, dan untuk
$q\in\Bbb Q$

\Centerline{$\{x:x\in E,\,f(x)\le q\}
=H\cap\{x:x\in G,\,f(x)\le q\}\in\Sigma_H$.}

\noindent Jadi $f\restr E$ bersifat $\Sigma_H$-terukur dan $f$ bersifat
$\mu_H$-terukur secara virtual.

\medskip

\quad{\bf (iii)} Andaikan sedikitnya salah satu integral terdefinisi. Maka
(ii) memberi tahu kita bahwa terdapat suatu himpunan $\mu$-koterabaikan
$E\in\Sigma$ sedemikian sehingga $\tilde f\restr E$ bersifat
$\Sigma$-terukur; dalam hal ini $f\restr H\cap E$ bersifat
$\Sigma_H$-terukur.

\medskip

\qquad\grheada\ Andaikan $f$ nonnegatif di seluruh domainnya. Maka
$\int_Hfd\mu_H$ dan $\int_X\tilde fd\mu$ keduanya terdefinisi dalam
$[0,\infty]$. Jika keduanya tak berhingga, kita dapat berhenti. Jika tidak,

\Centerline{$G=\{x:x\in E\cap H$, $f(x)<\infty\}
=\{x:x\in E$, $\tilde f(x)<\infty\}$}

\noindent harus bersifat koterabaikan. Tetapkan $g=f\restr G\cap H$; maka
$\tilde g=\tilde f\restr G$, sehingga
$g=f\,\,\text{hampir di mana-mana terhadap }\mu_H$ dan
$\tilde g=\tilde f\,\,\text{hampir di mana-mana terhadap }\mu$. Oleh karena itu
$\int_Hfd\mu_H=\int_Hg\,d\mu_H$ dan
$\int_X\tilde fd\mu=\int_X\tilde g\,d\mu$. Sekarang kita mengandaikan
sedikitnya salah satu dari keduanya berhingga. Namun dalam kasus ini kita
dapat menerapkan 131E untuk melihat bahwa
$\int_Hg\,d\mu=\int_X\tilde g\,d\mu$, sehingga
$\int_Hf\,d\mu=\int_X\tilde f\,d\mu$.

\medskip

\qquad\grheadb\ Secara umum, nyatakan $f$ sebagai $f^+-f^-$, dengan

\Centerline{$f^+(x)=\max(0,f(x))$,
\quad$f^-(x)=\max(0,-f(x))$}

\noindent untuk $x\in\dom f$. Maka $(f^+)\ssptilde=\tilde f^+$ dan
$(f^-)\ssptilde=\tilde f^-$. Jadi

\Centerline{$\int_Hfd\mu_H
=\int_Hf^+d\mu_H-\int_Hf^{-}d\mu_H
=\int_X\tilde f^+d\mu-\int_X\tilde f^{-}d\mu
=\int_X\tilde fd\mu$}

\noindent jika salah satu dari keempat ungkapan terdefinisi dalam
$[-\infty,\infty]$.

\medskip

{\bf (b)} Tetapkan $f=h\restr H$; maka
$(h\times\chi H)(x)=\tilde f(x)$ untuk setiap $x\in\dom h$, sehingga
(a-iii) memberi tahu kita bahwa

\Centerline{$\int_Xh\times\chi H\,d\mu=\int_X\tilde fd\mu
=\int_H(h\restr H)d\mu_H$}

\noindent jika salah satu dari ketiga integral terdefinisi dalam
$[-\infty,\infty]$.

\medskip

{\bf (c)} Dengan menetapkan $h^+(x)=\max(0,h(x))$ dan
$h^-(x)=\max(0,-h(x))$ untuk $x\in\dom h$, baik $\int_Xh^+d\mu$ maupun
$\int_Xh^-d\mu$ terdefinisi dalam $[0,\infty]$, dan paling banyak salah
satunya tak berhingga. Secara khusus, keduanya bersifat $\mu$-terukur secara
virtual dan terdefinisi $\mu$-hampir di mana-mana, sehingga hal yang sama
berlaku untuk $h^+\times\chi H$ dan $h^-\times\chi H$. Karena
$\int_Xh^+\times\chi Hd\mu\le\int_Xh^+d\mu$ dan
$\int_Xh^-\times\chi Hd\mu\le\int_Xh^-d\mu$, paling banyak salah satu dari
$\int_Xh^+\times\chi Hd\mu$, $\int_Xh^-\times\chi Hd\mu$ tak berhingga,
dan

\Centerline{$\int_Xh\times\chi Hd\mu
=\int_Xh^+\times\chi Hd\mu-\int_Xh^-\times\chi Hd\mu$}

\noindent terdefinisi dalam $[-\infty,\infty]$. Menurut (b) di atas,
$\int_H(h\restr H)d\mu_H$ terdefinisi dalam $[-\infty,\infty]$.

\medskip

{\bf (d)(i)} Andaikan $\int_Xg\,d\mu$ terdefinisi dalam
$[-\infty,\infty]$ dan
$h\times\chi H\le g\,\,\text{hampir di mana-mana terhadap }\mu$. Maka

\Centerline{$\int_H(g\restr H)d\mu_H=\int_Xg\times\chi Hd\mu$}

\noindent terdefinisi, menurut (c); dan karena $g(x)\ge 0$ untuk
$\mu$-hampir setiap $x\in X\setminus H$, berlaku
$g\times\chi H\leae g$. Jadi

\Centerline{$\overlineint_H(h\restr H)d\mu_H
\le\int_H(g\restr H)d\mu_H=\int_Xg\times\chi Hd\mu\le\int_Xg\,d\mu$.}

\noindent Karena $g$ sembarang,
$\overline{\int}_H(h\restr H)d\mu_H
\le\overline{\int}_Xh\times\chi H\,d\mu$.

\medskip

\quad{\bf (ii)} Andaikan $\int_Hfd\mu_H$ terdefinisi dalam
$[-\infty,\infty]$ dan
$h\restr H\le f\,\,\text{hampir di mana-mana terhadap }\mu_H$. Maka
$\int_X\tilde fd\mu$ terdefinisi dalam $[-\infty,\infty]$ dan
$h\times\chi H\le\tilde f\,\,\text{hampir di mana-mana terhadap }\mu$, sehingga

\Centerline{$\overlineint_Xh\times\chi H\,d\mu
\le\int_X\tilde fd\mu=\int_Hf\,d\mu_H$.}

\noindent Karena $f$ sembarang,
$\overline{\int}_Xh\times\chi H\,d\mu
\le\overline{\int}_H(h\restr H)d\mu_H$.
}%end of proof of 135I

\exercises{
\leader{135X}{Latihan dasar (a)} Kita mengatakan bahwa suatu himpunan
$G\subseteq[-\infty,\infty]$ bersifat {\bf terbuka} jika (i)
$G\cap \Bbb R$ terbuka dalam arti biasa sebagai subhimpunan $\Bbb R$
(ii) jika $\infty\in G$, maka terdapat suatu $a\in\Bbb R$ sedemikian
sehingga $\ocint{a,\infty}\subseteq G$ (iii) jika $-\infty\in G$ maka
terdapat suatu $a\in\Bbb R$ sedemikian sehingga
$\coint{-\infty,a}\subseteq G$. Tunjukkan bahwa keluarga $\frak T$ dari
subhimpunan-subhimpunan terbuka $[-\infty,\infty]$ mempunyai sifat-sifat
yang bersesuaian dengan (a)-(d) dari 1A2B.

\header{135Xb}{\bf (b)} Tunjukkan bahwa himpunan-himpunan Borel dari
$[-\infty,\infty]$ sebagaimana didefinisikan dalam 135C tepat merupakan
anggota-anggota aljabar-$\sigma$ atas subhimpunan-subhimpunan
$[-\infty,\infty]$ yang dibangkitkan oleh himpunan-himpunan terbuka
sebagaimana didefinisikan dalam 135Xa.

\sqheader 135Xc Definisikan $\phi:[-\infty,\infty]\to[-1,1]$ dengan
menetapkan

\Centerline{$\phi(-\infty)=-1$, \quad $\phi(x)=\tanh x
=\Bover{e^{2x}-1}{e^{2x}+1}$ jika $-\infty<x<\infty$, \quad
$\phi(\infty)=1$.}

\noindent Tunjukkan bahwa (i) $\phi$ merupakan isomorfisme urutan antara
$[-\infty,\infty]$ dan $[-1,1]$ (ii) untuk sembarang barisan
$\sequencen{u_n}$ dalam $[-\infty,\infty]$, $\sequencen{u_n}\to u$ jika
dan hanya jika $\sequencen{\phi(u_n)}\to\phi(u)$ (iii) untuk sembarang
himpunan $E\subseteq [-\infty,\infty]$, $E$ bersifat Borel dalam
$[-\infty,\infty]$ jika dan hanya jika $\phi[E]$ merupakan subhimpunan
Borel dari $\Bbb R$ (iv) suatu fungsi bernilai real $h$ yang terdefinisi
pada subhimpunan dari $[-\infty,\infty]$ bersifat terukur Borel jika dan hanya
jika $h\phi^{-1}$ bersifat terukur Borel.

\sqheader 135Xd Misalkan $X$ suatu himpunan, $\Sigma$ suatu
aljabar-$\sigma$ atas subhimpunan-subhimpunan $X$ dan $f$ suatu fungsi dari
subhimpunan dari $X$ ke $[-\infty,\infty]$. Tunjukkan bahwa $f$ terukur jika dan
hanya jika komposisi $\phi f$ terukur, dengan $\phi$ fungsi dari 135Xc.
Gunakan hal ini untuk mereduksi 135Ef dan 135Eg menjadi hasil-hasil yang
bersesuaian dalam \S121.

\header{135Xe}{\bf (e)} Misalkan $\phi:[-\infty,\infty]\to[-1,1]$ fungsi
yang dijelaskan dalam 135Xc. Tunjukkan bahwa fungsi-fungsi

\Centerline{$(t,u)\mapsto\phi(\phi^{-1}(t)+\phi^{-1}(u)):
[-1,1]^2\setminus\{(-1,1),(1,-1)\}\to[-1,1]$,}

\Centerline{$(t,u)\mapsto\phi(\phi^{-1}(t)\phi^{-1}(u)):
[-1,1]^2\to[-1,1]$,}

\Centerline{$(t,u)\mapsto\phi(\phi^{-1}(t)/\phi^{-1}(u)):
\bigl([-1,1]\times([-1,1]\setminus\{0\})\bigr)
\setminus\{(\pm 1,\pm 1)\}\to[-1,1]$}

\noindent bersifat terukur Borel. Gunakan hasil ini bersama 121K untuk
membuktikan 135Ee.

\header{135Xf}{\bf (f)} Dengan mengikuti konvensi-konvensi 135Ab dan 135Ad,
berikan uraian lengkap tentang kasus-kasus ketika
$uu'/vv'=(u/v)(u'/v')$ dan ketika $uw/vw=u/v$.

\spheader 135Xg Misalkan $(X,\Sigma,\mu)$ suatu ruang ukur dan andaikan
$E\in\Sigma$ mempunyai ukuran berhingga yang tak nol. Misalkan $f$ suatu
fungsi bernilai $[-\infty,\infty]$ yang terukur secara virtual dan
terdefinisi pada suatu subhimpunan dari $X$, dan andaikan $f(x)$ terdefinisi dan
lebih besar daripada $\alpha$ untuk hampir setiap $x\in E$. Tunjukkan bahwa
$\int_Ef>\alpha\mu E$.
%135F

\leader{135Y}{Latihan lanjutan (a)}
%\spheader 135Ya
Misalkan $X$ suatu himpunan dan $\Sigma$ suatu aljabar-$\sigma$ atas
subhimpunan-subhimpunan $X$. Tunjukkan bahwa jika $f:X\to[0,\infty]$
bersifat $\Sigma$-terukur, terdapat suatu barisan $\sequencen{E_n}$ dalam
$\Sigma$ sedemikian sehingga
$f=\sum_{n=0}^{\infty}\Bover1{n+1}\chi E_n$.

\spheader 135Yb Misalkan $(X,\Sigma,\mu)$ suatu ruang ukur, dan $f$, $g$
dua fungsi bernilai $[-\infty,\infty]$, yang masing-masing terdefinisi pada
suatu subhimpunan dari $X$, sedemikian sehingga $\int f$ dan $\int g$
keduanya terdefinisi dalam $[-\infty,\infty]$.
(i) Tunjukkan bahwa $\int f\vee g$ dan $\int f\wedge g$ terdefinisi dalam
$[-\infty,\infty]$, dengan $(f\vee g)(x)=\max(f(x),g(x))$,
$(f\wedge g)(x)=\min(f(x),g(x))$ untuk
$x\in\dom f\cap\dom g$. (ii) Tunjukkan bahwa
$\int f\vee g+\int f\wedge g=\int f+\int g$ dalam arti bahwa jika salah
satu jumlah terdefinisi dalam $[-\infty,\infty]$, maka yang lainnya juga
terdefinisi, dan keduanya sama.
%135F

\spheader 135Yc\dvAnew{2007}
Misalkan $(X,\Sigma,\mu)$ suatu ruang ukur,
$f:X\to[-\infty,\infty]$ suatu fungsi dan $g:X\to[0,\infty]$,
$h:X\to[0,\infty]$ fungsi-fungsi terukur. Tunjukkan bahwa
$\overline{\int}f\times(g+h)
=\overline{\int}f\times g+\overline{\int}f\times h$, dengan di sini kita
menafsirkan $\infty+(-\infty)$ sebagai $\infty$, seperti dalam 133L.
}%end of exercises

\endnotes{
\Notesheader{135} Saya menempatkan uraian ini dalam bagian tersendiri
sebagian karena panjangnya, dan sebagian karena saya ingin menekankan bahwa
teknik-teknik ini hanya bersifat sampingan terhadap gagasan-gagasan utama
jilid ini. Sesungguhnya yang saya coba lakukan di sini hanyalah memberikan
uraian yang koheren tentang bahasa yang lazim digunakan untuk menangani
berbagai kasus pinggiran. Sebagai aturan umum, `$\infty$' masuk ke dalam
argumen-argumen ini hanya sebagai singkatan untuk jenis-jenis hal yang
sepele. Ketika kita ingin memberikan nilai $\pm\infty$ kepada suatu fungsi,
entah hal itu terjadi pada suatu himpunan terabaikan -- dalam hal ini sering
kali tepat, meski sedikit kurang menenteramkan, untuk menganggap fungsi
tersebut tidak terdefinisi pada himpunan itu -- atau keadaan sudah benar-
benar lepas kendali dan teori hanya dapat memberi tahu kita sedikit hal yang
berguna.

Tentu saja tidak sulit memasukkan teori garis real diperluas langsung ke
dalam argumen-argumen Bab 12, sehingga hasil-hasil bagian ini menjadi
hasil-hasil dasar. Saya menghindari jalur ini sebagian sebagai upaya
mengurangi jumlah gagasan baru yang diperlukan dalam materi Bab 12 yang
secara teknis sangat menuntut -- karena saya percaya bahwa, terlepas dari
cara kita memperlakukan $\pm\infty$, mutlak perlu untuk dapat menangani
fungsi-fungsi yang terdefinisi secara parsial -- dan sebagian karena saya
tidak berpendapat bahwa garis real semestinya dianggap sebagai suatu
substruktur dari garis real diperluas. Menurut saya, keduanya merupakan
struktur yang berbeda dengan sifat-sifat yang berbeda, dan garis real asli
jauh lebih penting. Namun, dapat dikatakan bahwa dari segi gagasan-gagasan
yang dibahas dalam jilid ini keduanya begitu serupa sehingga ketika Anda
benar-benar akrab dengan karya ini Anda akan dapat berpindah dengan bebas
dari yang satu ke yang lain, bahkan sedemikian bebas sehingga Anda dapat
dengan aman membatasi pembedaan itu pada situasi-situasi formal,
misalnya ketika Anda menyajikan pernyataan suatu teorema.
}%end of notes

\discrpage
